\chapter{Conclusões}

{\color{red}
Escreva as conclusões aqui.
}


Backup Referencias (Alterar)

FAMA, E. F. Efficient capital markets: II. The Journal of Finance, v. 46, n. 5, p. 1575-1956, 1991. 

LO, A. W. The adaptive markets hypothesis: Market efficiency from an evolutionary perspective. Journal of Portfolio Management, v. 30, n. 5, p. 15-29, 2004.

PICASSO, A. et al. Technical analysis and sentiment embeddings for market trend prediction. Expert Systems with Applications, v. 135, n. 23, p. 60-70, 2019. 

YOO, P. D. et al. Machine learning techniques and use of event information for stock market prediction: A survey and evaluation. International Conference on Computational Intelligence for Modelling, Control and Automation and International Conference on Intelligent Agents, Web Technologies and Internet Commerce, v. 2, n. 1, p. 835–841, 2005.

BOUKTIF, S. et al. Augmented Textual Features-based Stock Market Prediction. IEEE Access, v. 8, n. 1, p. 40269-40282, 2020.
MIKOLOV, T. et al. Efficient Estimation of Word Representations in Vector Space. v. 1, n. 1, p. 1-12, 2013.

JIN, N. et al. Multi-Task Learning Model Based on Multi-Scale CNN and LSTM for Sentiment Classification. IEEE Access, v. 8, n. 1, p. 77060-77072, 2020.

GOMEZ-ADORNO, H. et al. A Convolutional Neural Network Approach for Gender and Language Variety Identification. Journal of Intelligent & Fuzzy Systems, v. 36, n. 5, p. 4845-4855, 2019.

KRIZHEVSKY, T. Web Scraping with Excel. IEEE Access, v. 8, n. 1, p. 40269-40282, 2020.

KIM, Y. Convolutional neural networks for sentence classification. Conference on Empirical Methods in Natural Language Processing (EMNLP), v. 1, n. 1, p. 1746–1751, 2014.

PANG, B. LEE, L. Seeing stars: Exploiting class relationships for sentiment categorization with respect to rating scales. 43rd Annu. Meeting Assoc. Comput. Linguistics (ACL), p. 115-124, 2005.

HOCHREITER, S. The vanishing gradient problem during learning recurrent neural nets and problem solutions. International Journal Uncertainty, Fuzziness Knowledge-Based Systems, v. 6, n. 2, p. 107-116, 1998.

ZHANG, S. et al., Multi-layer attention based CNN for target-dependent sentiment classification, Neural Process. Lett., v. 51, n. 3, p. 1–15, 2019.

WEN, S. et al., Memristive LSTM network for sentiment analysis. IEEE Transactions on Systems, Man, and Cybernetics, v. 51, n.3, p. 1794-1804, 2021.

Zhang, M. et al. Gated neural networks for targeted sentiment analysis. 13th AAAI Conference on Artificial Intelligence, v. 30, n. 1, p. 3087–3093, 2016.

MANIKANDAPRABHU, T. Web Scraping with Excel. International Scientific Journal of Engineering and Management, v. 2, n. 4, p. 1-3, 2023.
HESSEL, M. et al. Multi-task deep reinforcement learning with popart, AAAI Conference on Artificial Intelligence, v. 33, n. 1, p. 3796–3803, 2019.

LIU, P. QIU, X. HUANG, X. Deep multi-task learning with shared memory, Proceedings of the 2016 Conference on Empirical Methods in Natural Language Processing, p. 118-127, 2016.

YOUSIF, A. et al., ‘‘Multi-task learning model based on recurrent convolutional neural networks for citation sentiment and purpose classification, IEEE Access, v. 335, n. 1 p. 195–205, 2019.

LU, G. et al., Multi-task learning using variational auto-encoder for sentiment classification, Pattern Recognition Letters, v. 132, n. 1, p. 115–122, 2020.
